% science_template.tex
% See accompanying readme.txt for copyright statement, change log etc.

% Any modification of this template, including writing a paper using it,
% MUST rename the file i.e. use a different file name.

%%%%%%%%%%%%%%%% START OF PREAMBLE %%%%%%%%%%%%%%%

% Basic setup. Authors shouldn't need to adjust these commands.
% It's annoying, but please do NOT strip these into a separate file.
% They need to be included in this .tex for our production software to work.

% Use the basic LaTeX article class, 12pt text
\documentclass[12pt]{article}

% Science uses Times font. If you don't have this installed (most LaTeX installations will be
% fine) or prefer the old Computer Modern fonts, comment out the following line
\usepackage{newtxtext,newtxmath}
% Depending on your LaTeX fonts installation, you might get better results with one or both of these:
%\usepackage{mathptmx}
%\usepackage{txfonts}

% Allow external graphics files
\usepackage{graphicx}
\usepackage{microtype}

% Use US letter sized paper with 1 inch margins
\usepackage[letterpaper,margin=1in]{geometry}

% Double line spacing, including in captions
\linespread{1.5} % For some reason double spacing is 1.5, not 2.0!

% One space after each sentence
\frenchspacing

% Abstract formatting and spacing - no heading
\renewenvironment{abstract}
	{\quotation}
	{\endquotation}

% No date in the title section
\date{}

% Reference section heading
\renewcommand\refname{References and Notes}

% Figure and Table labels in bold
\makeatletter
\renewcommand{\fnum@figure}{\textbf{Figure \thefigure}}
\renewcommand{\fnum@table}{\textbf{Table \thetable}}
\makeatother

% Call the accompanying scicite.sty package.
% This formats citation numbers in Science style.
%\usepackage{scicite}

% Provides the \url command, and fixes a crash if URLs or DOIs contain underscores
\usepackage{url}

%%%%%%%%%%%% CUSTOM COMMANDS AND PACKAGES %%%%%%%%%%%%

% Authors can define simple custom commands e.g. as shortcuts to save on typing
% Use \newcommand (not \def) to avoid overwriting existing commands.
% Keep them as simple as possible and note the warning in the text below.

% Example:
\newcommand{\pcc}{\,cm$^{-3}$}	% per cm-cubed

% If your BibTeX uses shortcuts for journal names, define them here e.g.
% \newcommand{\apj}{Astrophys.~J.}

% Please DO NOT import additional external packages or .sty files.
% Those are unlikely to work with our conversion software and will cause problems later.
% Don't add any more \usepackage{} commands.


%%%%%%%%%%%%%%%% TITLE AND AUTHORS %%%%%%%%%%%%%%%%

% Title of the paper.
% Keep it short and understandable by any reader of Science.
% Avoid acronyms or jargon. Use sentence case.
\def\scititle{
	Decomposing gains in average wages across nativity status
}
% Store the title in a variable for reuse in the supplement (otherwise \maketitle deletes it)
\title{\bfseries \boldmath \scititle}

% Author and institution list.
% Institution numbers etc. should be hard-coded, do *not* use the \footnote command.
\author{
	% You can write out first names or use initials - either way is acceptable, but be consistent
	Patricia van Hissenhoven Flórez$^{1}$ and
	Xi Song$^{2}$ \and
	% Additional lines of authors should be inserted using the \and command (not \\)
	% Institution list, in a slightly smaller font
	\small$^{1}$Department of Sociology, University of Pennsylvania. \and\
	\small$^{2}$Department of Sociology, Columbia University.\and
	% Do not include street addresses, postal codes etc.
	%
	% Identify at least one corresponding author, with contact email address
	\small$^\ast$Corresponding author. Email: example@mail.com\and
	% Joint contributions can be indicated like this
	\small$^\dagger$These authors contributed equally to this work.
}

%%%%%%%%%%%%%%%%% END OF PREAMBLE %%%%%%%%%%%%%%%%


%%%%%%%%%%%%%%%% START OF MAIN TEXT %%%%%%%%%%%%%%%
\begin{document} 

% Insert the title and author list
\maketitle

% Abstract, in bold
% There are strict length limits, and not all formats have abstracts.
% Consult the journal instructions to authors for details.
% Do not cite any references in the abstract.
\begin{abstract} 
\end{abstract}

\section{Introduction}
This paper decomposes the differences in the growth of immigrant and native average wages between 2001 and 2023. It asks three research questions: 
\begin{enumerate}
    \item What are the trends in the occupational average wage growth between natives and foreign born?
    \item Are the changes in the gaps between natives and immigrants progressive or do the shifts occur at specific times?
    \item Are the gaps explained by changes in the occupational allocation of workers in occupations or within wage gains? Are these changes concentrated in few or many occupations?
    \item Which nativity group and occupations have benefited more from worker reallocation? Which has benefited more from wage gains?
\end{enumerate}



% The first paragraph of any Science paper does NOT have a heading
% Nor is it indented
\noindent

\section*{Data and Methods}

Leveraging the large sample size of the American Community Survey, we observe the over 2.6 million working age (25 to 65) respondents who received wages between 2001 and 2023. To assess time trends within occupation, the sample is further restricted to 416 occupations for which at least 1 worker is surveyed for every year (98.3 percent of observations and 91 percent of occupations). However, sample size for most occupations in most years is over 30. We stratify samples by nativity status, analyzing the growth in average wages for native and foreign born workers separately. All estimation of counts of workers in each occupation and occupational wages are weighted by the PERWT provided by the ACS. 

This paper takes an occupation-centered approach to understand broader job market trends. First, we provide descriptive graphs showing the trends in the average wages by nativity. All results are disaggregated by nativity status. Based on these results, we identify four periods of trends that we further disaggregate. These generally align with key events in the economy: 2001 to 2007, before the Great Recession; 2008 - 2011, the Great Recession; 2012- 2019, the recovery period; and 2020-2023, COVID-19 and its aftermath. 

The second set of analyses use a labor-market level decomposition to understand which interacted components that explain the difference in the wage gap over time. The goal is to discern the extent to which changes in occupational segregation versus diminishing (or increasing) wage gaps contribute to the decline in the gap. Note that this is a decomposition of the change in the wage gap, that is a decomposition of the difference -across time- of the difference in average wages by nativity group. 

 Equation 1 below illustrates that a nativity group's mean earnings are itself a weighted average of its occupation-specific earnings. At a given time $t$, for workers of nativity group $N$, this is the sum of the average wage of a given occupation $i$, $\bar{X_{Ni}}$, weighted by the proportion of workers of that nativity status $P_{Ni}$.  The gap in average occupational wage corresponds to the difference between the average occupational wage for natives minus foreign born workers.

Using the standard Kitagawa decomposition, we can decompose the native-immigrant earnings gap at time $t$, $\Delta_t$, into two effects. The first component of the decomposition is often known as the \emph{rate effect}, due to differences in group-specific rates (e.g., earnings). The second component is often known as the \emph{composition effect}, due to differences in group compositions (e.g., occupational structure).

\begin{align}
\Delta_t &= \text{mean earnings of natives} - \text{mean earnings of immigrants} \notag \\
&= \sum_i P_{Ni}^t \bar{X}_{Ni}^t - \sum_i P_{Ii}^t \bar{X}_{Ii}^t \\
&= \underbrace{\sum_i \left(\frac{P_{Ni}^t + P_{Ii}^t}{2}\right) \left(\bar{X}_{Ni}^t - \bar{X}_{Ii}^t\right)}_{\text{rate effect}} + \underbrace{\sum_i \left(\frac{\bar{X}_{Ni}^t + \bar{X}_{Ii}^t}{2}\right) \left(P_{Ni}^t - P_{Ii}^t\right)}_{\text{composition effect}}
\end{align}


To extend the Kitagawa decomposition to two time points, we further introduce four key concepts. First, the occupational share, corresponding to the proportion of workers within a given nativity group $N$ employed in occupation $i$. The average of the two proportions can be represented as $P_i^t = \frac{P_{i|N}^t + P_{i|I}^t}{2}$.  Second, the native-immigrant earnings gap within each occupation is $\delta_i^t = \bar{X}_{Ni}^t - \bar{X}_{Ii}^t$. Third, the occupation-level mean earnings, across all workers is $\bar{X}_i^t = \frac{\bar{X}_{Ni}^t + \bar{X}_{Ii}^t}{2}$. Here we consider a perfectly symmetric counterfactual for occupation $i$ by using the "simple" average of the mean earnings of native and immigrant workers. Doing so yields the neat two-term, fully additive decomposition. Fourth, we introduce the occupational segregation measure, the difference between the native and immigrant share within an occupation $S_i^t = P_{i|N}^t - P_{i|I}^t$.

Using these definitions, we can then rewrite the definition of $\Delta$ in equation (2) as follows:

\begin{align}
\Delta = \; & \sum_i \left(\frac{P_{i|N}^2 + P_{i|I}^2}{2}\right)\left(\bar{X}_{Ni}^2 - \bar{X}_{Ii}^2\right) + \sum_i \left(\frac{\bar{X}_{Ni}^2 + \bar{X}_{Ii}^2}{2}\right)\left(P_{i|N}^2 - P_{i|I}^2\right) \notag \\
& - \sum_i \left(\frac{P_{i|N}^1 + P_{i|I}^1}{2}\right)\left(\bar{X}_{Ni}^1 - \bar{X}_{Ii}^1\right) - \sum_i \left(\frac{\bar{X}_{Ni}^1 + \bar{X}_{Ii}^1}{2}\right)\left(P_{i|N}^1 - P_{i|I}^1\right)
\end{align}

Where, by plugging in the four terms and reorganizing them, yields a four-term decomposition. 


\begin{multline}
    \Delta = \underbrace{\sum_i \left(\frac{\bar{X}_i^2 + \bar{X}_i^1}{2}\right)\left(S_i^2 - S_i^1\right)}_{C1} + \underbrace{\sum_i \left(\frac{P_i^2 + P_i^1}{2}\right)\left(\delta_i^2 - \delta_i^1\right)}_{C2}\\
+ \underbrace{\sum_i \left(\frac{\delta_i^2 + \delta_i^1}{2}\right)\left(P_i^2 - P_i^1\right)}_{C4} + \underbrace{\sum_i \left(\frac{S_i^2 + S_i^1}{2}\right)\left(\bar{X}_i^2 - \bar{X}_i^1\right)}_{C3} 
\end{multline}

We can interpret the first component as the change in occupational nativity segregation and the second component as the within-occupation native-immigrant earnings gaps. The two latter components account for macro-structural changes across the labor market, regardless of the nativity status. The third component corresponds to the change in occupational earnings structure, accounting for the decline of jobs as they have become lower or higher paid in the job market and the fourth component shows the change in occupational composition as occupations  have become more on less common among all workers. 

While this decomposition provides insight into changes across the labor market, it makes it harder to discern whether natives or immigrants  benefited differentially from occupational growth. For example, if occupational segregation declined, it is difficult to discern whether the decline in segregation is occurring because of increased benefits for immigrants or declines in advantages for natives. For example, a decline in occupational segregation may occur if immigrants are entering higher paid and higher status jobs, but also if natives are entering lower paid and lower status jobs. To further discern this, we use a more conventional share-shift analysis to  decompose the growth of average occupational wages within each nativity group into three components. 

Recall that mean wages $\bar{X^N}_{N} = \sum_i P_{i, t}^N\bar{X}^N_{i, t}$, where $P^N_{i,N}$ denotes the share of workers of nativity $N$ in occupation $i$, at time $t$,  and $\bar{X}^N_{i, t}$ is the mean wage within occupation for that nativity group at that time. Changes in aggregate wages between periods $t=0$ and $t=1$ are decomposed into a between-occupation (reallocation) component, a within-occupation wage component, and a covariance term such that: 

\begin{equation}
	\Delta\bar{X} = \sum_i (P_{i,1} - P_{i,0})\bar{X}_{i,0} + \sum_{i,0}(\bar{X}_{i,1}- \bar{X}_{i,0}) + \sum_i (s_{s,1} - s_{o,0})(\bar{w}_{o,1}- \bar{w}_{o,0}) 
\end{equation}

Note that the $N$ subscript has been dropped because each of these estimations is done separately for each nativity group. From this we can understand which nativity group has leveraged occupational reallocation and wage gains differently. Furthermore, we can identify the specific occupations that have contributed more or less to the change. 




%%%%%%%%%%%%%%%% MAIN TEXT FIGURES %%%%%%%%%%%%%%%

\begin{figure} % Do NOT use \begin{figure*}
	\centering
	\includegraphics[width=0.6\textwidth]{example_figure} % for an image file named example_figure.*
	% Pick an appropriate width - in print, figures are usually one or two columns wide, which can
	% be approximated by 0.3\textwidth or 0.6\textwidth respectively. Use appropriate label sizes.

	% Captions go below figures
	\caption{\textbf{All captions must start with a short bold sentence, acting as a title.}
		Then explain what is being shown, the meanings of any line styles, plotting symbols etc.
		Multi-panel figures must label the panels A, B, C, etc. and refer to them in the caption
		like this: (\textbf{A}) Description of panel A. (\textbf{B}) Description of panel B.
		Captions are placed below figures.}
	\label{fig:example} % give each figure a logical label name
\end{figure}


%%%%%%%%%%%%%%%% MAIN TEXT TABLES %%%%%%%%%%%%%%%

\begin{table} % Do NOT use \begin{table*}
	\centering
	% Captions go above tables
	\caption{\textbf{All captions must start with a short bold sentence, acting as a title.}
		Then explain what is being listed in the table, the meaning of each column etc.
		Captions are placed above tables.}
	\label{tab:example} % give each table a logical label name
	
	\begin{tabular}{lccc} % four columns, alignment for each
		\\
		\hline
		Sample & $A$ & $B$ & $C$\\
		 & (unit) & (unit) & (unit)\\
		\hline
		First & 1 & 2 & 3\\
		Second & 4 & 6 & 8\\
		Third & 5 & 7 & 9\\
		\hline
	\end{tabular}
\end{table}



%%%%%%%%%%%%%%%% REFERENCES %%%%%%%%%%%%%%%

\clearpage % Clear all remaining figures and tables then start a new page

% The list of references goes after the main text and before the acknowledgements
% When preparing an initial submission, we recommend you use BibTeX, like this:
%
\bibliography{science_template} % for a file named science_template.bib
\bibliographystyle{sciencemag}

% After the paper has completed peer review and been revised ready for acceptance,
% you should comment out the lines above and copy-paste the contents of your .bbl
% file here instead. This will help ensure that our conversion software works correctly.
% Remember to re-run BibTeX first - check the timestamp!
%
% Example of the first three entries copy-pasted from science_template.bbl:
%
%\begin{thebibliography}{1}
%
%\bibitem{example}
%A.~N. {Author}, An example reference. \emph{Journal of Improbable Research}
%  \textbf{1}, 67 (2020).
%
%\bibitem{example2}
%F.~M. {Surname}, S.~{Author}, A second example. \emph{Interesting Research
%  Letters} \textbf{32}, 897 (2019).
%
%\bibitem{example_preprint}
%P.~{One}, P.~{Two}, P.~{Three}, {An unpublished preprint}. \emph{preprint}
%  (2021), arXiv:2101.12345.
%
%\end{thebibliography}


%%%%%%%%%%%%%%%% ACKNOWLEDGEMENTS %%%%%%%%%%%%%%%

\section*{Acknowledgments}
Here you can thank helpful colleagues who did not meet the journal's authorship criteria, or
provide other acknowledgements that don't fit the (compulsory) subheadings below.
Formatting requirements for each of these sections differ between the \textit{Science}-family
journals; consult the instructions to authors on the journal website for full details.
\paragraph*{Funding:}
List the grants, fellowships etc. that funded the research;
use initials to specify which author(s) were supported by each source.
Include grant numbers when appropriate or required by the funding agency.
For example: F.~A. was funded by the Generous Science Agency grant~2372.
\paragraph*{Author contributions:}
List each author’s contributions to the paper.
Use initials to abbreviate author names.
\paragraph*{Competing interests:}
Disclose any potential conflicts of interest for all authors, such as patent applications,
additional affiliations, consultancies, financial relationships etc.
See the journal editorial policies web page for types of competing interest that must be declared.
If there are no competing interests, state:
``There are no competing interests to declare.''
\paragraph*{Data, code and materials availability:}
Specify where the data, software, physical samples, simulation outputs or other materials
underlying the paper are archived.
These must be publicly accessible when the paper is published (without embargo) and enable
readers to fully reproduce all the results in the paper.
Contact the editor if you’re unsure what needs to be shared.

Our preference is for digital material to be deposited in a suitable non-profit online data or
software repository that assigns it a DOI.
Alternatively, an institutional repository, subject-based archive, commercial repository etc.
is acceptable, as are (short) supplementary tables or a machine-readable supplementary data file.
`Available on request’ or personal web pages are not allowed.

Cite the relevant DOI \cite{dataset}, URL \cite{example_url} or reference \cite{example2}
in this statement.
These \textit{do not} count towards the reference limit if they are only cited in the acknowledgements.
Be specific and state a unique identifier -- such as an accession number, software version number
or observation ID -- so readers can easily retrieve the exact material used.

Declare any restrictions on sharing or re-use -- such as a materials transfer agreement (MTA) or
legal restrictions -- which must be approved by your editor.
Unreasonable restrictions will preclude publication.
Consult the journal's editorial policies web page for more details.

If there are no physical materials, state:
``No physical materials were generated in this work.''


%%%%%%%%%%%%%%%% SUPPLEMENT LIST %%%%%%%%%%%%%%%

% List the contents of your Supplementary Materials, including the numbers of any
% supplementary figures, tables, external data files etc. and any references that are
% cited only in the supplement. In this example, refs. 7-8 are cited only in the supplement.
% Fill out your numbers accordingly and delete any lines that aren't applicable.
\subsection*{Supplementary materials}
Materials and Methods\\
Supplementary Text\\
Figs. S1 to S3\\
Tables S1 to S4\\
References \textit{(7-\arabic{enumiv})}\\ % automatically fills out the last reference number
% (filling out the other numbers automatically is possible but fiddly and liable to break)
Movie S1\\
Data S1

%%%%%%%%%%%%%%%% END OF MAIN TEXT %%%%%%%%%%%%%%%

\newpage

%%%%%%%%%%%%%%%% START OF SUPPLEMENT %%%%%%%%%%%%%%%

% Figures, tables, equations and pages in the supplement are numbered S1, S2 etc.
\renewcommand{\thefigure}{S\arabic{figure}}
\renewcommand{\thetable}{S\arabic{table}}
\renewcommand{\theequation}{S\arabic{equation}}
\renewcommand{\thepage}{S\arabic{page}}
\setcounter{figure}{0}
\setcounter{table}{0}
\setcounter{equation}{0}
\setcounter{page}{1} % not 0 as \newpage already started a supplementary page
% References continue the numbering from the main text.


%%%%%%%%%%%%%%%% SUPPLEMENT TITLE PAGE %%%%%%%%%%%%%%%

\begin{center}
\section*{Supplementary Materials for\\ \scititle}

% Author list for the supplement
% Indicate the corresponding authors, but do NOT include institutions here
% It would be nice if the template auto-generated this, but doing so is complicated...
First~Author$^{\ast\dagger}$,
A.~Scientist$^\dagger$,
Someone~E.~Else\\ % we're not in a \author{} environment this time, so use \\ for a new line
\small$^\ast$Corresponding author. Email: example@mail.com\\
\small$^\dagger$These authors contributed equally to this work.
\end{center}

% Fill out the numbers for each type of supplementary material,
% and delete any lines that aren't applicable.
% These are just example numbers that don't match the rest of this template.
\subsubsection*{This PDF file includes:}
Materials and Methods\\
Supplementary Text\\
Figures S1 to S3\\
Tables S1 to S4\\
Captions for Movies S1 to S2\\
Captions for Data S1 to S2

\subsubsection*{Other Supplementary Materials for this manuscript:}
Movies S1 to S2\\
Data S1 to S2

\newpage

%%%%%%%%%%%%%%%% MATERIALS AND METHODS %%%%%%%%%%%%%%%


\subsubsection*{Example supplement heading}

The two main sections of the supplement can be split up using headings.


%%%%%%%%%%%%%%%% SUPPLEMENTARY TEXT %%%%%%%%%%%%%%%

\subsection*{Methods Memo 1}
\subsubsection{Research question}
\begin{enumerate}
    \item Have the immigrant-native worker earnings gaps changed since 2000?
    \item If so, what mechanisms can explain the changes:
    \begin{itemize}
        \item \textbf{Mechanism 1 (sorting mechanism):} Has sorting into different occupations of immigrant workers changed over time? Check out (1) immigrant niches theory developed by Roger Waldinger at UCLA; (2) dual labor market theory by Portes and Massey; (3) occupational segregation by immigrant status.
        \item \textbf{Mechanism 2 (competition mechanism):} Has earnings competition between immigrant and native workers within the same occupation changed over time? Devaluation of immigrants' skills/education, discrimination.
    \end{itemize}
\end{enumerate}

\subsubsection*{Analytical Plan}

\begin{enumerate}
    \item Decompose immigrant-native earnings gaps into these mechanisms using mean earnings.
    \item Decompose immigrant-native earnings gaps into these mechanisms using the variances of earnings.
\end{enumerate}

\subsubsection*{Methods}

\subsection*{1. Decomposing Changes in Native-Immigrant Mean Earnings Gaps into Four Components}

We extend the Kitagawa (1955) decomposition to a two-group, two-time-points setting to examine how the earnings gap between immigrant and native workers changes over time. The decomposition consists of four components: (1) changes in occupational nativity segregation (how immigrants and natives are distributed across occupations); (2) changes in within-occupation nativity gap in earnings (how much more immigrants or natives earn within the same occupation); (3) changes in the overall occupation-level earnings (which occupations are high- or low-paying overall); and (4) changes in the overall occupational composition (relative sizes of occupations regardless of nativity).

First, we define the native-immigrant earnings gap at time $t$ as

\begin{equation}
\Delta_t = \bar{X}_N^t - \bar{X}_I^t
\end{equation}

where $\bar{X}_N^t$ and $\bar{X}_I^t$ refer to the average earnings of native and immigrant workers, respectively.

For two time points (1 and 2), our outcome of interest is

\begin{equation}
\Delta = \Delta_2 - \Delta_1
\end{equation}

Using the standard Kitagawa decomposition, we can decompose the native-immigrant earnings gap at time $t$, $\Delta_t$, into two effects:

\begin{align}
\Delta_t &= \text{mean earnings of natives} - \text{mean earnings of immigrants} \notag \\
&= \sum_i P_{Ni}^t \bar{X}_{Ni}^t - \sum_i P_{Ii}^t \bar{X}_{Ii}^t \\
&= \underbrace{\sum_i \left(\frac{P_{Ni}^t + P_{Ii}^t}{2}\right) \left(\bar{X}_{Ni}^t - \bar{X}_{Ii}^t\right)}_{\text{rate effect}} + \underbrace{\sum_i \left(\frac{\bar{X}_{Ni}^t + \bar{X}_{Ii}^t}{2}\right) \left(P_{Ni}^t - P_{Ii}^t\right)}_{\text{composition effect}}
\end{align}

The equation illustrates that each nativity group's mean earnings are itself a weighted average of its occupation-specific earnings. The term $P_{Ni}^t$ (and $P_{Ii}^t$) refers to occupational shares of native (and immigrant) workers. The first component of the decomposition is often known as the \emph{rate effect}, due to differences in group-specific rates (e.g., earnings). The second component is often known as the \emph{composition effect}, due to differences in group compositions (e.g., occupational structure).

To extend the Kitagawa decomposition to two time points, we further introduce the following concepts.

\paragraph{1. Occupational share:}

\begin{align}
P_{i|N}^t &= \frac{N_i^t}{N^t} \\
P_{i|I}^t &= \frac{I_i^t}{I^t} \\
P_i^t &= \frac{P_{i|N}^t + P_{i|I}^t}{2}
\end{align}

Here $N_i^t$ (and $I_i^t$) refers to the number of native (and immigrant) workers in occupation $i$ at $t$; $N^t$ (and $I^t$) refers to the total number of native (and immigrant) workers at $t$. We average the two groups' shares $P_{i|N}^t$, $P_{i|I}^t$ to get $P_i^t$.

\paragraph{2. Native-immigrant earnings gap within each occupation:}

\begin{equation}
\delta_i^t = \bar{X}_{Ni}^t - \bar{X}_{Ii}^t
\end{equation}

where $\bar{X}_{Ni}^t$ (and $\bar{X}_{Ii}^t$) is the mean earnings of native (and immigrant) workers in occupation $i$ at $t$.

\paragraph{3. Occupation-level mean earnings (all workers):}

\begin{equation}
\bar{X}_i^t = \frac{\bar{X}_{Ni}^t + \bar{X}_{Ii}^t}{2}
\end{equation}

Here we consider a perfectly symmetric counterfactual for occupation $i$ in equation (4) by using the ``simple'' average of the mean earnings of native and immigrant workers. Doing so yields the neat two-term, fully additive decomposition in equation (4).\footnote{Alternatively, we can use a more complicated weighted average as $\bar{X}_i^t = \dfrac{N_i^t \bar{X}_{Ni}^t + I_i^t \bar{X}_{Ii}^t}{N_i^t + I_i^t}$, which is each occupation's actual workforce-weighted mean earnings. This would result in a more complicated decomposition with interaction terms.}

\paragraph{4. Occupational segregation measure (difference in native vs. immigrant share within occupation $i$):}

\begin{equation}
S_i^t = P_{i|N}^t - P_{i|I}^t
\end{equation}

Using these definitions, we can then rewrite the definition of $\Delta$ in equation (4) as follows:

\begin{align}
\Delta = \; & \sum_i \left(\frac{P_{i|N}^2 + P_{i|I}^2}{2}\right)\left(\bar{X}_{Ni}^2 - \bar{X}_{Ii}^2\right) + \sum_i \left(\frac{\bar{X}_{Ni}^2 + \bar{X}_{Ii}^2}{2}\right)\left(P_{i|N}^2 - P_{i|I}^2\right) \notag \\
& - \sum_i \left(\frac{P_{i|N}^1 + P_{i|I}^1}{2}\right)\left(\bar{X}_{Ni}^1 - \bar{X}_{Ii}^1\right) - \sum_i \left(\frac{\bar{X}_{Ni}^1 + \bar{X}_{Ii}^1}{2}\right)\left(P_{i|N}^1 - P_{i|I}^1\right)
\end{align}

Plugging in the definitions of $P_i^t$, $delta_i^t$, $\bar{X}_i^t$, and $S_i^t$,

\begin{equation}
\Delta = \sum_i P_i^2 delta_i^2 + \sum_i \bar{X}_i^2 S_i^2 - \sum_i P_i^1 delta_i^1 - \sum_i \bar{X}_i^1 S_i^1
\end{equation}

Reorganizing the terms,

\begin{equation}
\Delta = \sum_i \left(P_i^2 \delta_i^2 - P_i^1 \delta_i^1\right) + \sum_i \left(\bar{X}_i^2 S_i^2 - \bar{X}_i^1 S_i^1\right)
\end{equation}

Applying the Kitagawa decomposition to $P_i^2 \delta_i^2 - P_i^1 \delta_i^1$ and $\bar{X}_i^2 S_i^2 - \bar{X}_i^1 S_i^1$, respectively, we get

\begin{equation}
\Delta = \underbrace{\sum_i \left(\frac{P_i^2 + P_i^1}{2}\right)\left(\delta_i^2 - \delta_i^1\right)}_{C2} + \underbrace{\sum_i \left(\frac{\delta_i^2 + \delta_i^1}{2}\right)\left(P_i^2 - P_i^1\right)}_{C4} + \underbrace{\sum_i \left(\frac{S_i^2 + S_i^1}{2}\right)\left(\bar{X}_i^2 - \bar{X}_i^1\right)}_{C3} + \underbrace{\sum_i \left(\frac{\bar{X}_i^2 + \bar{X}_i^1}{2}\right)\left(S_i^2 - S_i^1\right)}_{C1}
\end{equation}

We can interpret each component as follows.

\begin{description}
    \item[C1: Change in occupational nativity segregation] $\displaystyle\sum_i \left(\frac{\bar{X}_i^2 + \bar{X}_i^1}{2}\right)\left(S_i^2 - S_i^1\right)$ shows how shifts in the native vs. immigrant worker mix within occupations (the $S_i$'s) matter, weighted by each occupation's mean earnings.

    \item[C2: Change in within-occupation native-immigrant earnings gaps] $\displaystyle\sum_i \left(\frac{P_i^2 + P_i^1}{2}\right)\left(\delta_i^2 - \delta_i^1\right)$ shows how much of the nativity gap change owes to widening or narrowing of gaps within each occupation, weighted by the average occupational share.

    \item[C3: Change in occupational earnings structure] $\displaystyle\sum_i \left(\frac{S_i^2 + S_i^1}{2}\right)\left(\bar{X}_i^2 - \bar{X}_i^1\right)$ shows how changes in the average earnings of each occupation (for all workers) affect the gap, scaled by that occupation's average nativity imbalance.

    \item[C4: Change in occupational composition] $\displaystyle\sum_i \left(\frac{\delta_i^2 + \delta_i^1}{2}\right)\left(P_i^2 - P_i^1\right)$ shows how the overall growth or shrinkage of entire occupations (regardless of nativity) shifts the gap, given the within-occupation earnings gaps.
\end{description}

\end{document}

%%%%%%%%%%%%%%%% SUPPLEMENTARY FIGURES %%%%%%%%%%%%%%%

\begin{figure} % Do not use \begin{figure*}
	\centering
	\includegraphics[width=0.6\textwidth]{example_figure} % for an image file named example_figure.*
	% Pick an appriopriate width for the size of the image

	% Captions go below figures
	\caption{\textbf{All captions must start with a short bold sentence, acting as a title.}
		Follow the same style as main text figures.
		If the design is substantially the same as another figure, avoid repeating information
		e.g. say `Same as Figure~\ref{fig:example}, but for the control sample.'}
	\label{fig:sup_example} % give each figure a logical label name
\end{figure}

%%%%%%%%%%%%%%%% SUPPLEMENTARY TABLES %%%%%%%%%%%%%%%

\begin{table} % Do not use \begin{table*}
	\centering
	% Captions go above tables
	\caption{\textbf{All captions must start with a short bold sentence, acting as a title.}
		Follow the same style as main text tables.
		If the design is similar to previous tables, avoid repetition by refering back to them.}
	\label{tab:sup_example} % give each table a logical label name

	\begin{tabular}{lccr} % four columns, alignment for each
		\\
		\hline
		A & B & C & D\\
		\hline
		1 & 2 & 3 & 4\\
		2 & 4 & 6 & 8\\
		3 & 5 & 7 & 9\\
		\hline
	\end{tabular}
\end{table}


%%%%%%%%%%% CAPTIONS FOR OTHER SUPPLEMENTARY FILES %%%%%%%%%%

\clearpage % Clear all remaining figures and tables then start a new page

\paragraph{Caption for Movie S1.}
\textbf{All captions must start with a short bold sentence, acting as a title.}
Then explain what is shown in the supplementary video file.
Give as much detail as you would for a figure e.g. explain axes, color maps etc.
If the video is an animated equivalent of one of the static figures, state e.g.
`Animated version of Figure~\ref{fig:example}.'

\paragraph{Caption for Data S1.}
\textbf{All captions must start with a short bold sentence, acting as a title.}
Then explain what is included in the supplementary data file.
Give as much detail as you would for a table e.g. explain the meaning of every column,
units used, any special notation etc.


%%%%%%%%%%%%%%%% SUPPLEMENTARY REFERENCES %%%%%%%%%%%%%%%

% Do NOT include a reference list in the supplement.
% All references must be in a single list at the end of the main text.
% The copyeditors will ensure that the correct reference list appears with each version of the paper
% (print, HTML, PDF, mobile app, metadata for bibliographic databases etc.)

\end{document}
% End of science_template.tex

